\section{An optimal test for symmetric powers of an elliptic curve}
\label{sec:sym}
Fix a non-CM elliptic curve $E/\mathbb Q$ and an integer $d\ge1$.
Write $L_j(s)=L(s,\mathrm{Sym}^j E)$ in unitary normalization
(the center is $1/2$), and put $L_0=\zeta$.
For each good prime $p$, let $2\cos\theta_p=(p+1-\#E(\mathbb F_p))/\sqrt p$
with $0\le\theta_p\le\pi$. Define $V_0=1$, $V_1=z$ and
$V_{j+1}=zV_j-V_{j-1}$. Thus $V_j(2\cos\theta)=
\sin((j+1)\theta)/\sin\theta$, with endpoint values given by continuity.

\paragraph{An extremal polynomial.}
For a nonzero polynomial $f\ge0$ on $[-2,2]$ of degree at most $2d$,
write $f=\sum_{j=0}^{2d} a_j V_j$. The largest possible ratio $a_1/a_0$ is
\[
 u=2\cos\frac{\pi}{d+2}.
\]
The unique optimizer up to positive scale is
\[
 f_d(z)=\left(\sum_{j=0}^d V_j(u)V_j(z)\right)^2 .
\]
Every coefficient $a_0,\ldots,a_{2d}$ of this optimizer is strictly positive.
The ratio measures the coefficient of the ordinary trace when the mean
of the test is fixed. Other objectives, such as minimizing sample variance,
give different optimization problems.

\emph{Proof.}
The $V_j$ are orthonormal for
$d\mu(z)=\sqrt{4-z^2}\,dz/(2\pi)$ by sine orthogonality.
Write $\theta_i=i\pi/(d+2)$, $z_i=2\cos\theta_i$ and
$w_i=2\sin^2\theta_i/(d+2)$, for $1\le i\le d+1$.
The rule $\int g\,d\mu=\sum_i w_i g(z_i)$ holds through degree $2d+1$.
Indeed, divide $g$ by $V_{d+1}$, which vanishes at all nodes. The quotient
has degree at most $d$, so the integral of the multiple of $V_{d+1}$
vanishes by orthogonality. Discrete sine orthogonality gives the formula for
the remainder. Now $a_1/a_0=\int zf\,d\mu/\int f\,d\mu$ is a weighted
average of the nodes, so it is at most $z_1=u$.
Equality forces $f$ to have a double zero at every other node, fixing its
shape up to scale. Discrete sine orthogonality gives
$\sum_{j=0}^d V_j(z_1)V_j(z_i)=0$ for $i>1$, so $f_d$ attains equality.
Each $V_j(u)$ is positive for $0\le j\le d$. The product rule
$V_j V_k=V_{j+k}+V_{j+k-2}+\cdots+V_{|j-k|}$ follows from the sine formula.
For each $0\le m\le2d$, some $j,k\le d$ have $j+k=m$, so $a_m>0$.
\hfill$\square$

\paragraph{A fourth-power criterion.}
Use the extremizer just obtained and form the positive sum
\[
 S_f(x)=\sum_{\substack{p^k\le x,\ k\ge1\\p\ {\rm good}}}
 (\log p)\, f(2\cos(k\theta_p))\left(\frac1{p^k}-\frac1x\right).
\]
Write $L_j^{(S)}$ for the Euler product with bad primes omitted and set
\[
 D(s)=-\sum_{j=0}^{2d}a_j\frac{(L_j^{(S)})'}{L_j^{(S)}}(s),\qquad
 C=\lim_{s\to1}\left(D(s)-\frac{a_0}{s-1}\right)-a_0,\qquad
 R_f(x)=S_f(x)-a_0\log x-C .
\]
Bad Euler factors have normalized inverse roots of modulus at most one,
so their removal introduces no zeros in $\Re s>0$. For fixed $j\ge1$,
cuspidal automorphy, entireness of the completed function and functional
equation are available by modularity and Newton--Thorne, Theorem A and
Corollary B (\emph{Publ. Math. IH\'ES} 134, 2021,
\href{https://doi.org/10.1007/s10240-021-00126-4}{doi:10.1007/s10240-021-00126-4}).
We use Jacquet and Shalika's nonvanishing theorem on $\Re s=1$
(\emph{A non-vanishing theorem for zeta functions of $\mathrm{GL}_n$},
\emph{Invent. Math.} 38 (1976), 1--16). In particular the displayed limit exists.

\textbf{Theorem.}
All nontrivial zeros of $L_0,\ldots,L_{2d}$ have real part $1/2$
if and only if, for each $\epsilon>0$, there exists $B_\epsilon>0$ with
\[
 R_f(n^4)\le B_\epsilon n^{-2+4\epsilon}
 \qquad\hbox{for all sufficiently large integers }n.
\]
Here $E$ and $d$ are fixed. The theorem yields a single positive prime
sum for a finite family of conjectures. Under the premise of the question,
the $j=0$ member is already supplied.

\emph{Interpolation.}
In $t=1/x$, the function $S_f(1/t)$ is convex. Each summand is a
nonnegative coefficient times $(p^{-k}-t)_+$. Between
$t_-=(n+1)^{-4}$ and $t_+=n^{-4}$ it is bounded by its secant.
The secant of $-a_0\log t$ exceeds that function by at most
$a_0(t_+-t_-)^2/(8t_-^2)=O_d(n^{-2})$, by the interpolation
remainder and the second derivative $a_0/t^2$. Since $x\asymp n^4$,
the sample bounds imply $R_f(x)\le B'_\epsilon x^{-1/2+\epsilon}$.
This argument accounts for the fourth powers and the sign of the test.

\emph{Zeros from the upper bound.}
In $\Re s>1$ termwise integration gives
\[
 H(s):=\int_1^\infty R_f(x)x^{-s}\,dx
  =\frac{D(s)}{s(s-1)}-\frac{a_0}{(s-1)^2}-\frac C{s-1}.
\]
Indeed $V_j(2\cos(k\theta_p))$ is the trace of the $k$th power of
the normalized Frobenius in $\mathrm{Sym}^j$. The definition of $C$
removes the singularity at $1$. At a zero $\rho$ in the open critical strip,
the residue is
\[
 \frac{\sum_j a_j\,\mathrm{ord}_{\rho}L_j}{\rho(1-\rho)}.
\]
It is positive at every real zero and nonzero at every nonreal zero,
including a zero shared by several factors.
For $\beta=1/2+\epsilon<1$, take the eventually nonnegative function
$B'_\epsilon x^{\beta-1}-R_f(x)$. Changing it on a compact interval
affects its transform by an entire function. If the Mellin integral of
this function has abscissa of convergence $\sigma_c>\beta$, Landau's
theorem forces a singularity there. The continued transform is meromorphic
with at most simple poles there, and a real pole has negative residue
by the displayed formula. This is impossible for an integral of a
nonnegative function as $s\downarrow\sigma_c$.
Consequently $H$ is holomorphic for $\Re s>\beta$. Letting
$\epsilon\downarrow0$ and using the functional equations proves the claim.
This argument allows arbitrary orders at the central point.

\emph{The reverse implication.}
Assume the stated zero location and fix $0<\epsilon<1/2$.
Invert $D(s)/(s(s-1))$ on $\Re s=2$ against $x^{s-1}$.
The inverse of each prime-power term is $(p^{-k}-x^{-1})_+$.
Shift the line to $\Re s=1/2+\epsilon$. The double pole at $1$
contributes $a_0\log x+C$. For clarity, the usual justification of
the shift uses the fixed-degree unitary $L$-functions' order-one Hadamard
products, their gamma factors and the local zero count
$N(T+1)-N(T)=O_{E,d}(\log(2+T))$.
Subtracting the logarithmic derivative at $2+it$ from that at
$\sigma+it$ gives an absolutely convergent zero sum: the zeros with
$|\Im\rho-t|\le1$ contribute $O_{E,d,\epsilon}(\log(2+|t|))$ since
$|\sigma+it-\rho|\ge\epsilon$. More distant zeros contribute the same
bound by summing unit bands with weights $O((1+|t-\Im\rho|)^{-2})$.
Stirling's formula bounds the gamma terms. Thus $D(\sigma+it)$ is
$O_{E,d,\epsilon}(\log(2+|t|))$ uniformly on the shifted strip for
$|t|\ge2$. The quadratic denominator gives absolute convergence on the
new line and vanishing horizontal integrals. This gives the stronger
two-sided bound $|R_f(x)|\ll_{E,d,\epsilon} x^{-1/2+\epsilon}$.

\paragraph{The first four symmetric powers.}
For $d=2$, $u=\sqrt2$ and $f_2(z)=(z^2+\sqrt2 z)^2$.
Its character coefficients are $(4,4\sqrt2,5,2\sqrt2,1)$.
This gives an explicit target for an arithmetic upper bound at $x=n^4$.
The determinant estimates in the repository count congruence classes and
quadratic characters. A corresponding divisibility or height estimate
with these elliptic Frobenius weights is an additional arithmetic input.
The extremal problem and the sparse criterion are solved here. The bound
for this curve, and hence symmetric-power GRH, remains open in this analysis.

For the curve $y^2+y=x^3-x$, the discriminant is $37$ and
$j=48^3/37$ is nonintegral, which excludes CM. The accompanying
sample calculation counts points at every good prime up to $10^4$
and includes each admissible prime power. It gives
$S_f(16)=1.17616326$ and $S_f(10^4)=19.38460159$.
The table also gives consecutive differences of $S_f(x)-4\log x$,
in which the unknown additive constant cancels. These values illustrate
the observable at finite $x$.

\paragraph{How well can congruence weights approximate this test?}
For any fixed modulus $q$ and real function $g$ on its unit classes, write
$\pi_S(X)$ for the number of good primes $p\le X$ coprime to $q$.
The squared error for the local observable in the preceding example obeys
\[
 \lim_{X\to\infty}\frac{1}{\pi_S(X)}
 \sum_{\substack{p\le X\ {\rm good}\\p\nmid q}}
 \bigl(f_2(2\cos\theta_p)-g(p\bmod q)\bigr)^2
 =
 66+\frac1{\varphi(q)}\sum_{b\in(\mathbb Z/q\mathbb Z)^\times}
 (4-g(b))^2 .
\]
Indeed, in each unit residue class the average of
$V_j(2\cos\theta_p)$ tends to zero for $1\le j\le8$. Expand the class
indicator in Dirichlet characters and apply the prime number theorem to
the standard $L$-functions of $\mathrm{Sym}^j E$ twisted by these characters.
Each twist is cuspidal of degree $j+1>1$. Its standard $L$-function is
holomorphic and nonzero on $\Re s=1$, and its normalized prime-power traces
are bounded by $j+1$. These facts give the prime number theorem for the
trace by the usual Tauberian argument (add a sufficiently large constant
multiple of the zeta logarithmic derivative to bound the real and imaginary
coefficients by nonnegative ones, then subtract the ordinary prime number
theorem). Prime powers of exponent at least two cost $o(X)$ in the
log-weighted sum. Partial summation removes the logarithm.
Consequently the required average on class $b$ is its semicircle integral.
Orthonormality gives the squared fluctuation
$(4\sqrt2)^2+5^2+(2\sqrt2)^2+1^2=66$.
This quantifies precisely which local information survives every fixed
congruence test. It leaves open cancellation in a smoothed sum and the use
of moduli growing with the cutoff.
